\chapter{Introdução}
\label{chap:introducao}

Os sistemas operativos modernos distribuem carga entre núcleos de \ac{CPU} recorrendo a heurísticas clássicas~\cite{LinuxCFS}
que, apesar de eficazes em produção, tomam decisões subótimas em cenários de elevada contenção de recursos.

Em paralelo, métodos de otimização quântica, como o \ac{QAOA}~\cite{Farhi2014QAOA},
têm sido investigados em problemas combinatoriais que podem ser formulados como
\ac{QUBO}~\cite{Kochenberger2014QUBO,Glover2019QUBO}.
 
No entanto, a era \ac{NISQ} impõe limitações de escala e ruído que
tornam difícil avaliar quando um escalonador quântico oferece vantagens concretas sobre alternativas
clássicas~\cite{Preskill2018}. Preskill antecipa que esta barreira será ultrapassada com hardware da ordem
do \textit{megaquop} --- máquinas capazes de executar $10^6$ operações quânticas com erro
tolerável~\cite{Preskill2025} --- fornecendo um referencial concreto para situar a relevância futura
deste tipo de abordagem.

Este relatório apresenta uma plataforma híbrida quântico-clássica para testar o problema
de atribuição de processos a núcleos de \acs{CPU}. A plataforma não substitui o
escalonador do kernel; isola a decisão de atribuição processo-núcleo,
modela-a como um \acs{QUBO}, resolve-a com \acs{QAOA} via PennyLane~\cite{Bergholm2018PennyLane},
e valida os resultados por comparação com um \textit{solver} exato de \textit{brute force}
quando o tamanho da instância o permite. O objetivo não é demonstrar superioridade
quântica, mas caracterizar o comportamento do \acs{QAOA} neste domínio em função da
escala do problema.

\section{Formulação do Problema}
\label{sec:formulacao-problema}

Dado um conjunto de processos com utilização de \acs{CPU} conhecida e um número fixo de núcleos, o sistema deve
atribuir cada processo a exatamente um núcleo. O objetivo de otimização é o equilíbrio de carga: a carga
total de cada núcleo deve aproximar-se o mais possível da carga média global. A solução deve ainda respeitar
restrições de atribuição exclusiva --- nenhum processo pode ficar por atribuir ou ser atribuído a múltiplos
núcleos simultaneamente.

Expressar estas restrições e este objetivo numa formulação \acs{QUBO}, compatível com otimizadores quânticos,
constitui o problema central deste projeto.
